\sect{माघ-शुक्ल-जया-एकादशी-माहात्म्यम्}
\label{sec:padma-magha-shukla-jaya}

\ganapatyadiDhyanam
\padmaPuranam

\dnsub{अथ पञ्चचत्वारिंशोऽध्यायः॥४५}

\uvacha{युधिष्ठिर उवाच}

\twolineshloka
{साधु कृष्ण त्वया प्रोक्तमादिदेवो भवान् प्रभो}
{स्वेदजा अण्डजाश्चैव उद्भिज्जाश्च जरायुजाः}% ॥१॥

\twolineshloka
{तेषां कर्ता विकर्ता त्वं पालकः क्षयकारकः}
{माघस्य कृष्णपक्षे तु षट्तिला कथिता त्वया}% ॥२॥

\twolineshloka
{शुक्ले च का भवेद् देव कथयस्व प्रसादतः}
{किं नाम को विधिस्तस्याः को देवस्तत्र पूज्यते}% ॥३॥

\uvacha{श्रीकृष्ण उवाच}

\twolineshloka
{कथयिष्यामि राजेन्द्र शुक्ले माघस्य या भवेत्}
{जया नामेति विख्याता सर्वपापहरा परा}% ॥४॥

\twolineshloka
{पवित्रा पापहन्त्री च कामदा मोक्षदा नृणाम्}
{ब्रह्महत्यापहन्त्री च पिशाचत्वविनाशिनी}% ॥५॥

\twolineshloka
{नैव तस्या व्रते चीर्णे प्रेतत्वं जायते नृणाम्}
{नातः परतरा काचित् पापघ्नी मोक्षदायिनी}% ॥६॥

\twolineshloka
{एतस्मात् कारणाद् राजन् कर्तव्या सा प्रयत्नतः}
{श्रूयतां राजशार्दूल कथा पौराणिकी शुभा}% ॥७॥

\twolineshloka
{पङ्कजे च पुराणेऽस्या महिमा कथितो मया}
{एकदा नाकलोके वै इन्द्रो राज्यं चकार ह}% ॥८॥

\twolineshloka
{देवास्तत्र सुखेनैव निवसन्ति मनोरमे}
{पीयूषपाननिरता अप्सरोगणसेविताः}% ॥९॥

\twolineshloka
{नन्दनं तु वनं तत्र पारिजातोपसेवितम्}
{रमयन्ति रमन्तेऽत्र अप्सरोभिर्दिवौकसः}% ॥१०॥

\twolineshloka
{एकदा रममाणोऽसौ देवेन्द्रः स्वेच्छया नृप}
{नर्तयामास वै हर्षात् पञ्चाशत्कोटिनायकः}% ॥११॥

\twolineshloka
{गन्धर्वास्तत्र गायन्ति गन्धर्वः पुष्पदन्तकः}
{चित्रसेनस्तु तत्रैव चित्रसेनसुता तथा}% ॥१२॥

\twolineshloka
{मालिनीति च नाम्ना तु चित्रसेनस्य योषिता}
{मालिन्यास्तु समुत्पन्ना पुष्पदन्ती च नामतः}% ॥१३॥

\twolineshloka
{पुष्पदन्तस्य पुत्रोऽसौ माल्यवान्नाम नामतः}
{पुष्पदन्त्याश्च रूपेण माल्यवानतिमोहितः}% ॥१४॥

\twolineshloka
{तया ह्येवं कटाक्षैश्च माल्यवांश्च वशीकृतः}
{लावण्यं रूपसम्पन्नं तस्या रूपं निशामय}% ॥१५॥

\twolineshloka
{बाहू तस्याश्च कामेन कण्ठपाशौ कृताविव}
{कर्णायते तु नयने रक्तान्ते घूर्णिते तथा}% ॥१६॥

\twolineshloka
{कर्णौ तु शोभनौ तस्याः कुण्डलाभ्यां नृपोत्तम}
{कम्बुग्रीवायुता सैव दिव्याभरणभूषिता}% ॥१७॥

\twolineshloka
{पीनोन्नतौ कुचौ तस्यास्तौ हेमकलशाविव}
{मध्यं क्षामं च चार्वङ्ग्या मुष्टिग्राह्यमनुत्तमम्}% ॥१८॥

\twolineshloka
{नितम्बौ विस्तृतौ चास्या विस्तीर्णा जघनस्थली}
{चरणौ शोभमानौ च रक्तोत्पलसमद्युती}% ॥१९॥

\twolineshloka
{ईदृश्या पुष्पवत्या स माल्यवानतिमोहितः}
{शक्रस्य परितोषाय नृत्यार्थं तौ समागतौ}% ॥२०॥

\twolineshloka
{गायमानौ तु तौ तत्र अप्सरोगणसेवितौ}
{मदनाभिपरीताङ्गौ पुष्पदन्ती च माल्यवान्}% ॥२१॥

\twolineshloka
{परस्परानुरागेण व्यामोहवशमागतौ}
{न शुद्धगानं गायेतां चित्तभ्रमसमन्वितौ}% ॥२२॥

\twolineshloka
{बद्धदृष्टी तथाऽन्योन्यं कामबाणवशं गतौ}
{ज्ञात्वा लेखर्षभस्तत्र सङ्गतं मानसं तयोः}% ॥२३॥

\twolineshloka
{तालक्रियामानलोपात् तथा गीतविसर्जनात्}
{चिन्तयित्वा तु मघवा ह्यवमानं तथाऽऽत्मनः}% ॥२४॥

\twolineshloka
{कुपितश्च तयोरर्थे शापं दास्यन्निदं जगौ}
{धिग्धिग्वां पतितौ मूढावाज्ञाभङ्ग कृतौ मम}% ॥२५॥

\twolineshloka
{युवां पिशाचौ भवतां दम्पतीभावधारिणौ}
{मर्त्यलोकमनुप्राप्तौ भुञ्जानौ कर्मणः फलम्}% ॥२६॥

\twolineshloka
{एवं मघवता शप्तावुभौ दुःखितमानसौ}
{हिमवन्तं गिरिं प्राप्ताविन्द्रशापाद्विमोहितौ}% ॥२७॥

\twolineshloka
{उभौ पिशाचतां प्राप्तौ दारुणं दुःखमेव च}
{सन्तप्तमानसौ तत्र हिमकृच्छ्रगतावुभौ}% ॥२८॥

\twolineshloka
{गन्धर्वत्वमप्सरस्त्वं न जानीतो विमौहितौ}
{पीड्यमानौ निदाघेन देहपातकजेन च}% ॥२९॥

\twolineshloka
{न निशायां सुखं शान्तिं लभेते कर्मपीडितौ}
{परस्परं वादमानौ चेरतुर्गिरिगह्वरे}% ॥३०॥

\twolineshloka
{पीड्यमानौ तु शीतेन तुषारप्रभवेन तौ}
{दन्तघर्षं प्रकुर्वाणौ रोमाञ्चितवपुर्धरौ}% ॥३१॥

\twolineshloka
{ऊचे पिशाचः स तदा तां पत्नीं स्वां पिशाचिकाम्}
{किमनल्पकृतं पापं दारुणं रोमहर्षणम्}% ॥३२॥

\twolineshloka
{येन प्राप्तं पिशाचत्वं स्वेन दुष्कृतकर्मणा}
{नरकं दारुणं मत्वा पिशाचत्वं च दुःखदम्}% ॥३३॥

\twolineshloka
{तस्मात् सर्वप्रयत्नेन पातकं न समाचरेत्}
{इति चिन्तापरौ तत्र तावास्तां दुःखकर्षितौ}% ॥३४॥

\twolineshloka
{दैवयोगात्तयोः प्राप्ता माघस्यैकादशी तिथिः}
{जयानामेति विख्याता तिथीनामुत्तमा तिथिः}% ॥३५॥

\twolineshloka
{तस्मिन् दिने तु सम्प्राप्ते तावाहारविवर्जितौ}
{आसाते तत्र नृपते जलपानविवर्जितौ}% ॥३६॥

\twolineshloka
{न कृतो जीवघातश्च न पत्रफलभक्षणम्}
{अश्वत्थस्य समीपे तौ सर्वदा दुःखसंयुतौ}% ॥३७॥

\twolineshloka
{रविरस्तङ्गतो राजंस्तथैव स्थितयोस्तयोः}
{प्राप्ता चैव निशा घोरा दारुणा प्राणहारिणी}% ॥३८॥

\twolineshloka
{वेपमानौ ततस्तौ तु ततः सुषुपतः क्षितौ}
{परस्परेण संलग्नौ गात्रयोर्भुजयोरपि}% ॥३९॥

\twolineshloka
{न निद्रा न रतं तत्र न तौ सौख्यमविन्दताम्}
{एवं तौ राजशार्दूल शापेनेन्द्रस्य पीडितौ}% ॥४०॥

\twolineshloka
{इत्थं तयोर्दुःखितयोर्निर्जगाम निशीथिनी}
{मार्तण्ड उदयं प्राप्तो द्वादशी दिवसागमे}% ॥४१॥

\twolineshloka
{मया तु राजशार्दूल तयोर्मुक्तिर्धृता हृदि}
{जयायाः सुव्रतं चीर्णं रात्रौ जागरणं कृतम्}% ॥४२॥

\twolineshloka
{तस्माद् व्रतप्रभावाच्च यथाजातं तथा शृणु}
{द्वादशीदिवसे प्राप्ते तथा चीर्णे जया व्रते}% ॥४३॥

\twolineshloka
{विष्णोः प्रभावान्नृपते पिशाचत्वं तयोर्गतम्}
{पुष्पदन्ती माल्यवतौ पूर्वरूपौ बभूवतुः}% ॥४४॥

\twolineshloka
{पुरातनस्नेहयुतौ  पूर्वालङ्कारधारिणौ}
{विमानमधिरूढौ तौ गतौ नाके मनोरमे}% ॥४५॥

\twolineshloka
{देवेन्द्रस्याग्रतो गत्वा प्रणामं चक्रतुर्मुदा}
{तथाविधौ तु तौ दृष्ट्वा मघवा विस्मितोऽब्रवीत्}% ॥४६॥

\uvacha{इन्द्र उवाच}

\twolineshloka
{वद तं केन पुण्येन पिशाचत्वं हि वां गतौ}
{मम शापं च सम्प्राप्तौ केन देवेन मोचितौ}% ॥४७॥

\uvacha{माल्यवानुवाच}

\twolineshloka
{वासुदेवप्रसादेन जयायास्तु व्रतेन च}
{पिशाचत्वं गतं स्वामिंस्तव भक्तिप्रभावतः}% ॥४८॥

\twolineshloka
{इति श्रुत्वा तु मघवा प्रत्युवाच पुनस्तथा}
{पवित्रौ पावनौ जातौ वन्दनीयौ ममापि च}% ॥४९॥

\twolineshloka
{हरिवासरकर्तारौ विष्णुभक्तिपरायणौ}
{हरिवासरलीना ये ये च कृष्णपरायणाः}% ॥५०॥

\twolineshloka
{अस्माकमपि मर्त्यास्ते पूज्याश्चैव न संशयः}
{विहरस्व यथासौख्यं पुष्पदन्त्या सुरालये}% ॥५१॥

\uvacha{कृष्ण उवाच}

\twolineshloka
{एतस्मात् कारणाद् राजन् कर्तव्यो हरिवासरः}
{जया तु राजशार्दूल ब्रह्महत्यापहारिणी}% ॥५२॥

\twolineshloka
{सर्वदानानि तेनैव सर्वयज्ञान्यशेषतः}
{दत्तानि कारिताश्चैव जयायास्तु व्रतं कृतम्}% ॥५३॥

\twolineshloka
{कल्पकोटिं वसेत् तावद् वैकुण्ठे मोदते ध्रुवम्}
{पठनाच्छ्रवणाद् राजन्नग्निष्टोमफलं लभेत्}% ॥५४॥

॥इति श्रीपाद्मे महापुराणे पञ्चपञ्चाशत्साहस्र्यां संहितायामुत्तरखण्डे उमापतिनारदसंवादान्तर्गतकृष्णयुधिष्ठिरसंवादे माघ-शुक्ल-जया-एकादशी-माहात्म्यं नाम पञ्चचत्वारिंशोऽध्यायः॥४५॥

\genMangalaShloka

\hyperref[sec:ekadashi_mahatmyam_padma_puranam]{\closesub}
\clearpage

